\documentclass[11pt,a4paper]{report}

\usepackage[utf8]{inputenc}
\usepackage[T1]{fontenc}
\usepackage[french]{babel}
\usepackage{lmodern}
\usepackage{geometry}
\usepackage{graphicx}
\usepackage{hyperref}
\usepackage{xcolor}
\usepackage{booktabs}
\usepackage{longtable}
\usepackage{array}
\usepackage{enumitem}
\usepackage{listings}
\usepackage{tikz}
\usepackage{fancyhdr}
\usepackage{titlesec}

\geometry{margin=2.5cm}
\setlength{\headheight}{14pt}
\hypersetup{
  colorlinks=true,
  linkcolor=blue!60!black,
  urlcolor=blue!70!black,
  citecolor=green!50!black
}

\definecolor{codebg}{RGB}{245,245,250}
\definecolor{codeframe}{RGB}{200,200,220}

\lstset{
  backgroundcolor=\color{codebg},
  frame=single,
  rulecolor=\color{codeframe},
  basicstyle=\ttfamily\small,
  breaklines=true,
  columns=fullflexible,
  showstringspaces=false,
  literate={é}{{\'e}}1 {è}{{\`e}}1 {ê}{{\^e}}1 {à}{{\`a}}1
           {ù}{{\`u}}1 {ô}{{\^o}}1 {û}{{\^u}}1 {ç}{{\c{c}}}1
           {É}{{\'E}}1 {è}{{\`e}}1 {à}{{\`a}}1 {—}{---}1
}

\pagestyle{fancy}
\fancyhf{}
\fancyhead[L]{\textsc{Target OS --- Documentation CRM}}
\fancyhead[R]{\thepage}
\renewcommand{\headrulewidth}{0.4pt}

\titleformat{\chapter}{\Large\bfseries\color{blue!70!black}}{\thechapter}{1em}{}
\titleformat{\section}{\large\bfseries}{\thesection}{1em}{}
\titleformat{\subsection}{\normalsize\bfseries}{\thesubsection}{1em}{}

\begin{document}

% ---------------------------------------------------------------------------
\title{
  \vspace{2cm}
  {\Huge\bfseries Target OS}\\[0.5cm]
  {\LARGE Documentation technique du CRM}\\[0.3cm]
  {\large Plateforme de prospection B2B assistée par IA}\\[1.5cm]
  {\small \today}
}
\author{Target OS Dashboard}
\date{}
\maketitle
\thispagestyle{empty}
\tableofcontents
\newpage

% ---------------------------------------------------------------------------
\chapter{Vue d'ensemble}

\section{Présentation du projet}

\textbf{Target OS} est un CRM (Customer Relationship Management) orienté prospection B2B, conçu pour permettre aux \emph{prospecteurs} commerciaux de valider, prioriser et engager des leads identifiés automatiquement par l'intelligence artificielle.

Le système repose sur une architecture \textbf{Next.js} (App Router) côté frontend, avec \textbf{Supabase} comme base de données PostgreSQL hébergée, et des workflows \textbf{n8n} pour l'enrichissement automatique des prospects et la génération de rapports d'audit brand de type \emph{Awwards}.

\section{Objectifs métier}

\begin{itemize}[leftmargin=*]
  \item Centraliser les prospects générés par l'IA dans un pipeline visualisable.
  \item Permettre au prospecteur de \textbf{valider}, \textbf{approuver} et \textbf{contacter} chaque lead.
  \item Générer et partager des \textbf{rapports d'audit brand} personnalisés via un lien public (\texttt{/audit/\{slug\}}).
  \item Déclencher un chat proactif Tawk.to lorsque le prospect consulte les faiblesses de son entreprise (Mode Sniper).
  \item Suivre les performances du prospecteur via une section statistiques dédiée.
  \item Prendre des notes libres sur chaque lead pour le suivi terrain.
\end{itemize}

\section{Stack technique}

\begin{table}[h]
\centering
\begin{tabular}{@{}ll@{}}
\toprule
\textbf{Couche} & \textbf{Technologie} \\
\midrule
Framework frontend & Next.js 16.3.2 (App Router) \\
UI & React 19.2.8, Tailwind CSS 4, shadcn/ui \\
Composants UI & @base-ui/react \\
Icônes & lucide-react \\
Langage & TypeScript 5 \\
Base de données & Supabase (PostgreSQL) \\
Client BDD & @supabase/supabase-js 2.112.4 \\
Automatisation & n8n (externe) \\
Chat live & Tawk.to (embed JavaScript) \\
Typographie & Plus Jakarta Sans, Geist Mono \\
\bottomrule
\end{tabular}
\caption{Stack technique du projet \texttt{target-os-dashboard}}
\end{table}

\section{Structure du dépôt}

\begin{lstlisting}
target-os-dashboard/
|-- src/
|   |-- app/                    # Routes Next.js (App Router)
|   |   |-- page.tsx            # Dashboard pipeline (/)
|   |   |-- prospecteur/        # Stats prospecteur (/prospecteur)
|   |   |-- prospects/[id]/     # Fiche prospect interne
|   |   |-- audit/[slug]/       # Rapport public client
|   |   `-- api/prospects/      # API REST internes
|   |-- components/             # Composants React
|   |-- hooks/                  # Hooks personnalisés
|   |-- lib/                    # Utilitaires et logique métier
|   `-- types/                  # Types TypeScript (schéma canonique)
|-- supabase/                   # Migrations SQL
|-- n8n/                        # Script n8n pour rapports HTML
|-- docs/                       # Documentation
`-- .env.local                  # Variables d'environnement
\end{lstlisting}

% ---------------------------------------------------------------------------
\chapter{Modèle de données}

\section{Table \texttt{public.prospects}}

La table \texttt{prospects} est la source de vérité unique du CRM. Toutes les entités métier (lead, entreprise, audit, suivi) sont stockées dans cette table.

\subsection{Schéma complet}

\begin{longtable}{@{}p{3.5cm}p{2cm}p{1.5cm}p{7cm}@{}}
\toprule
\textbf{Colonne} & \textbf{Type} & \textbf{Requis} & \textbf{Description} \\
\midrule
\endhead
\texttt{id} & UUID & Oui & Clé primaire (auto-générée) \\
\texttt{created\_at} & TIMESTAMPTZ & Oui & Date de création (défaut : NOW()) \\
\texttt{slug} & TEXT & Non & Identifiant URL unique pour le rapport public \\
\midrule
\texttt{prenom} & TEXT & Non & Prénom du décisionnaire \\
\texttt{nom} & TEXT & Non & Nom du décisionnaire \\
\texttt{email} & TEXT & Oui & Email de contact \\
\texttt{poste} & TEXT & Non & Fonction / poste occupé \\
\midrule
\texttt{entreprise} & TEXT & Oui & Nom de l'entreprise cible \\
\texttt{url} & TEXT & Non & Site web de l'entreprise \\
\texttt{secteur} & TEXT & Non & Secteur d'activité \\
\texttt{taille\_entreprise} & TEXT & Non & Taille (effectifs, etc.) \\
\texttt{chiffre\_affaires} & TEXT & Non & Chiffre d'affaires estimé \\
\texttt{annee\_creation} & TEXT & Non & Année de fondation \\
\midrule
\texttt{ia\_score} & INTEGER & Non & Score de qualification IA (0--100) \\
\texttt{analyse\_site} & TEXT & Non & Analyse textuelle du site web \\
\texttt{forces} & TEXT & Non & Forces identifiées par l'IA \\
\texttt{faiblesses} & TEXT & Non & Faiblesses / points de rupture \\
\texttt{proposition\_commerciale} & TEXT & Non & Proposition commerciale générée \\
\texttt{script\_email} & TEXT & Non & Email de prospection prêt à l'envoi \\
\texttt{html\_rapport} & TEXT & Non & Rapport HTML brand complet (n8n) \\
\midrule
\texttt{statut} & TEXT & Oui & Statut pipeline (défaut : \emph{À valider}) \\
\texttt{notes} & TEXT & Non & Notes libres du prospecteur \\
\bottomrule
\caption{Schéma canonique de la table \texttt{public.prospects}}
\end{longtable}

\subsection{Index et contraintes}

\begin{itemize}[leftmargin=*]
  \item \textbf{Index} sur \texttt{statut}, \texttt{ia\_score} (DESC), \texttt{entreprise}.
  \item \textbf{Index unique partiel} sur \texttt{slug} (WHERE slug IS NOT NULL).
  \item \textbf{RLS activé} avec policies ouvertes pour \texttt{anon} et \texttt{authenticated} (SELECT, INSERT, UPDATE, DELETE).
\end{itemize}

\subsection{Valeurs de statut}

Les statuts sont des chaînes libres. Le code applique une logique de filtrage par mots-clés :

\begin{itemize}[leftmargin=*]
  \item \textbf{À valider} --- contient \texttt{valider} (statut par défaut).
  \item \textbf{Approuvé / Envoyé} --- contient \texttt{approuv} ou \texttt{envoy}.
  \item Toute autre valeur est classée dans \emph{Autres}.
\end{itemize}

\subsection{Constantes TypeScript}

Le fichier \texttt{src/types/database.types.ts} définit les groupes de champs :

\begin{lstlisting}
PROSPECT_IDENTITY_FIELDS   // prenom, nom, email, poste
PROSPECT_COMPANY_FIELDS    // entreprise, url, secteur, ...
PROSPECT_AI_CORE_FIELDS    // ia_score, analyse_site, forces, ...
PROSPECT_AI_FIELDS         // + html_rapport
PROSPECT_TRACKING_FIELDS   // statut, notes
\end{lstlisting}

Et les requêtes SELECT optimisées :
\begin{itemize}[leftmargin=*]
  \item \texttt{PROSPECT\_LIST\_SELECT} --- vue liste dashboard (sans \texttt{html\_rapport}).
  \item \texttt{PROSPECT\_DETAIL\_SELECT} --- fiche prospect (sans \texttt{html\_rapport}, chargé à la demande).
  \item \texttt{PROSPECT\_REPORT\_SELECT} --- uniquement \texttt{html\_rapport}.
\end{itemize}

% ---------------------------------------------------------------------------
\chapter{Routes et pages}

\section{Pages publiques et internes}

\begin{table}[h]
\centering
\begin{tabular}{@{}llp{7cm}@{}}
\toprule
\textbf{Route} & \textbf{Type} & \textbf{Description} \\
\midrule
\texttt{/} & Interne & Dashboard pipeline : KPIs + grille de cartes prospects \\
\texttt{/prospecteur} & Interne & Section stats du prospecteur \\
\texttt{/prospects/[id]} & Interne & Fiche détail d'un prospect \\
\texttt{/audit/[slug]} & Public & Rapport brand / audit client (sans auth) \\
\bottomrule
\end{tabular}
\end{table}

Toutes les pages de données utilisent \texttt{dynamic = "force-dynamic"} (pas de cache statique).

\subsection{Page Pipeline (\texttt{/})}

\begin{itemize}[leftmargin=*]
  \item Hero «~Vos prospects, sous contrôle~».
  \item 4 cartes KPI : total, à valider, approuvés, score moyen IA.
  \item Grille de cartes (\texttt{ProspectsBoard}) avec recherche textuelle et filtres par statut.
  \item Chaque carte permet : ouvrir la fiche, copier le lien audit, approuver le lead.
  \item Tri par \texttt{ia\_score} décroissant.
\end{itemize}

\subsection{Page Prospecteur (\texttt{/prospecteur})}

Section dédiée dans la sidebar pour les statistiques du commercial :

\begin{itemize}[leftmargin=*]
  \item Leads chauds (score $\geq$ 75) et froids ($<$ 50).
  \item Taux d'approbation (\%).
  \item Nombre de liens clients actifs (prospects avec \texttt{slug}).
  \item Notes saisies par le prospecteur.
  \item Activité récente (ajouts semaine / mois).
  \item Répartition par statut (barres de progression).
  \item Top 5 prospects par score IA (liens cliquables).
  \item Nom prospecteur par défaut : \emph{Prospecteur} (placeholder multi-utilisateurs futur).
\end{itemize}

\subsection{Fiche prospect (\texttt{/prospects/[id]})}

Interface «~Mission Control~» pour un lead individuel :

\paragraph{Hero}
Initiales entreprise, nom, badges statut/score, pills contact, bandeau meta (taille, site, CA, date).

\paragraph{Onglet Dossier}
\begin{enumerate}[leftmargin=*]
  \item \textbf{Notes prospecteur} --- textarea éditable, sauvegarde optimiste via API.
  \item \textbf{Identité \& entreprise} --- grille bento des champs d'enrichissement.
  \item \textbf{Audit IA} --- navigation verticale : Analyse, Forces, Faiblesses, Proposition.
  \item \textbf{Script email} --- bloc terminal dark, copie en un clic.
\end{enumerate}

\paragraph{Onglet Rapport}
Iframe embarqué de \texttt{/audit/\{slug\}}, préchargé en arrière-plan pour performance.

\paragraph{Sidebar sticky}
Score IA avec jauge, actions rapides (copier email, visiter site, audit client), résumé fiche.

\subsection{Audit public (\texttt{/audit/[slug]})}

Page sans chrome dashboard, accessible au client final :

\begin{enumerate}[leftmargin=*]
  \item Recherche du prospect par \texttt{slug} en base.
  \item Si slug inconnu $\rightarrow$ 404.
  \item Si \texttt{html\_rapport} vide $\rightarrow$ page \emph{En attente} (n8n n'a pas encore généré le rapport).
  \item Si rapport présent $\rightarrow$ rendu HTML via \texttt{dangerouslySetInnerHTML} (scripts strippés).
  \item Intégration Tawk.to + Mode Sniper (voir chapitre~\ref{ch:sniper}).
  \item Métadonnées : \texttt{robots: noindex, nofollow}.
\end{enumerate}

% ---------------------------------------------------------------------------
\chapter{API REST interne}

Quatre endpoints sous \texttt{/api/prospects/[id]/} :

\begin{longtable}{@{}llp{3cm}p{6cm}@{}}
\toprule
\textbf{Méthode} & \textbf{Route} & \textbf{Body} & \textbf{Action} \\
\midrule
\endhead
GET & \path{/report} & --- & Retourne \texttt{html\_rapport} uniquement (lazy load) \\
PATCH & \path{/statut} & \texttt{\{statut\}} & Met à jour le statut pipeline \\
PATCH & \path{/notes} & \texttt{\{notes\}} & Sauvegarde notes prospecteur (trim, null si vide) \\
POST & \path{/slug} & --- & Génère un slug unique depuis le nom entreprise \\
\bottomrule
\caption{Endpoints API internes}
\end{longtable}

Tous utilisent \texttt{createSupabaseClient()} avec la clé anon. Erreurs JSON \texttt{\{ error: string \}}.

\subsection{Génération de slug}

Algorithme dans \texttt{POST /api/prospects/[id]/slug} :
\begin{enumerate}[leftmargin=*]
  \item Si slug existant $\rightarrow$ retour immédiat.
  \item Sinon : \texttt{slugify(entreprise)} (minuscules, tirets, sans accents).
  \item En cas de collision : suffixe numérique (\texttt{entreprise-2}, \texttt{entreprise-3}, \ldots).
  \item Fallback : \texttt{entreprise-\{id[0:8]\}}.
\end{enumerate}

URL publique : \texttt{\{NEXT\_PUBLIC\_APP\_URL\}/audit/\{slug\}}.

% ---------------------------------------------------------------------------
\chapter{Composants principaux}

\section{Navigation et shell}

\begin{description}[style=nextline, leftmargin=3cm]
  \item[\texttt{DashboardShell}] Layout global : sidebar 72px, header sticky, mesh gradient background.
  \item[Sidebar] 3 entrées : Pipeline (/), Prospecteur (/prospecteur), Radar IA (désactivé).
  \item[Header] Label contextuel selon la route + badge «~Système actif~».
\end{description}

\section{Pipeline}

\begin{description}[style=nextline, leftmargin=3cm]
  \item[\texttt{DashboardKpi}] 4 cartes glass avec icônes et dégradés.
  \item[\texttt{ProspectsBoard}] Recherche + filtres pills + grille responsive.
  \item[\texttt{ProspectCard}] Carte lead : anneau score SVG, initiales, statut, actions.
  \item[\texttt{AuditLinkActions}] Génération/copie/prévisualisation lien client (card ou inline).
\end{description}

\section{Fiche prospect}

\begin{description}[style=nextline, leftmargin=3cm]
  \item[\texttt{ProspectDetailHero}] Hero éditorial avec grille de fond.
  \item[\texttt{ProspectDetailIdentity}] Grille bento identité/entreprise.
  \item[\texttt{ProspectDetailAudit}] Navigation latérale + panneaux colorés par type d'audit.
  \item[\texttt{ProspectDetailSidebar}] Score, jauge, actions, résumé.
  \item[\texttt{ProspectNotes}] Notes éditables avec feedback optimiste.
  \item[\texttt{ProspectScoreRing}] Anneau SVG + badge Lead chaud / Potentiel / À qualifier.
  \item[\texttt{AuditReportEmbed}] Iframe vers \texttt{/audit/slug} avec prefetch.
\end{description}

\section{Audit public}

\begin{description}[style=nextline, leftmargin=3cm]
  \item[\texttt{PublicHtmlReport}] Rendu HTML sécurisé + Tawk + sniper.
  \item[\texttt{TawkToWidget}] Chargement script embed Tawk.to.
  \item[\texttt{SniperProactiveBubble}] Bulle message personnalisé post-déclenchement.
  \item[\texttt{AuditPending}] Page d'attente si rapport non généré.
\end{description}

\section{Stats prospecteur}

\begin{description}[style=nextline, leftmargin=3cm]
  \item[\texttt{ProspectorStatsPanel}] Panel complet : KPIs étendus, activité, répartition, top 5.
\end{description}

% ---------------------------------------------------------------------------
\chapter{Bibliothèque utilitaire (\texttt{src/lib/})}

\begin{table}[h]
\centering
\small
\begin{tabular}{@{}lp{9cm}@{}}
\toprule
\textbf{Fichier} & \textbf{Rôle} \\
\midrule
\texttt{supabase.ts} & Factory client Supabase (anon key) \\
\texttt{get-prospects.ts} & Fetch liste prospects triée par score \\
\texttt{dashboard-stats.ts} & Calcul KPIs + stats prospecteur \\
\texttt{prospect-stats.ts} & Stats individuelles prospect (préparé, non branché UI) \\
\texttt{prospect-utils.ts} & Formatage, badges statut, noms complets \\
\texttt{slug.ts} & \texttt{slugify()}, \texttt{getPublicAuditUrl()} \\
\texttt{html-report.ts} & Strip scripts HTML, préparation rendu public \\
\texttt{sniper-message.ts} & Message sniper + événement CustomEvent \\
\texttt{utils.ts} & \texttt{cn()} --- merge classes Tailwind \\
\bottomrule
\end{tabular}
\end{table}

% ---------------------------------------------------------------------------
\chapter{Intégration Tawk.to et Mode Sniper}\label{ch:sniper}

\section{Tawk.to}

Widget de chat live intégré uniquement sur les pages d'audit public :

\begin{itemize}[leftmargin=*]
  \item Script : \texttt{https://embed.tawk.to/\{PROPERTY\_ID\}/\{WIDGET\_ID\}}.
  \item Variables : \texttt{NEXT\_PUBLIC\_TAWK\_PROPERTY\_ID}, \texttt{NEXT\_PUBLIC\_TAWK\_WIDGET\_ID}.
  \item Attributs visiteur : nom (prénom ou entreprise) + email.
  \item API utilisée : \texttt{setAttributes}, \texttt{showWidget}, \texttt{hideWidget}, \texttt{maximize}, \texttt{addEvent}.
\end{itemize}

\section{Mode Sniper (déclenchement comportemental)}

Hook \texttt{useSmartTawkTrigger} sur la page audit public :

\begin{enumerate}[leftmargin=*]
  \item \textbf{Observation} --- \texttt{IntersectionObserver} sur la section faiblesses du rapport HTML.
  \item \textbf{Sélecteurs cibles} --- \texttt{.mini-card.negative}, \texttt{[data-audit-section='faiblesses']}, ou titres contenant «~faiblesse~».
  \item \textbf{Délai} --- 15 secondes à 60\% de visibilité.
  \item \textbf{Action} --- ouverture Tawk (\texttt{maximize}) + événement \texttt{target\_os\_sniper\_faiblesses}.
  \item \textbf{Message} --- \emph{«~Re-bonjour \{Prénom\} ! Tu as vu le point rouge concernant les faiblesses de \{Entreprise\} ? On règle ça quand ?~»}
  \item \textbf{Bulle UI} --- \texttt{SniperProactiveBubble} affiche le message personnalisé (l'API Tawk ne permet pas d'injecter un message agent côté client).
\end{enumerate}

Le template HTML n8n inclut \texttt{class="mini-card negative" data-audit-section="faiblesses"} sur la section faiblesses pour servir de cible au trigger.

% ---------------------------------------------------------------------------
\chapter{Intégration n8n}

\section{Rôle de n8n}

n8n est le moteur d'automatisation externe responsable de :

\begin{enumerate}[leftmargin=*]
  \item \textbf{Collecte} --- trigger Google Sheets ou autre source de leads.
  \item \textbf{Enrichissement IA} --- analyse site, forces, faiblesses, proposition, score.
  \item \textbf{Insertion Supabase} --- création/mise à jour des lignes \texttt{prospects}.
  \item \textbf{Génération rapport} --- production du HTML brand + slug via script Code node.
\end{enumerate}

\section{Script \texttt{n8n/generate-brand-report.js}}

Ce script n8n :
\begin{itemize}[leftmargin=*]
  \item Parse la réponse JSON de l'IA (\texttt{ia\_score}, \texttt{analyse\_site}, \texttt{forces}, \texttt{faiblesses}, \texttt{proposition\_commerciale}).
  \item Lit les données Google Sheets (Entreprise, Prénom, Nom, Poste, URL, Secteur).
  \item Génère un rapport HTML dark style Awwards avec section faiblesses instrumentée pour le Mode Sniper.
  \item Produit en sortie : \texttt{\{ slug, html\_rapport \}} pour insertion Supabase.
\end{itemize}

\section{Champs attendus de n8n}

Pour un fonctionnement optimal, n8n doit envoyer à chaque insert/update :
\begin{itemize}[leftmargin=*]
  \item Tous les champs identité et entreprise.
  \item Les champs IA (\texttt{ia\_score}, \texttt{analyse\_site}, \texttt{forces}, \texttt{faiblesses}, \texttt{proposition\_commerciale}, \texttt{script\_email}).
  \item \texttt{slug} --- identifiant URL unique.
  \item \texttt{html\_rapport} --- rapport HTML complet.
\end{itemize}

% ---------------------------------------------------------------------------
\chapter{Migrations SQL Supabase}

\begin{table}[h]
\centering
\small
\begin{tabular}{@{}lp{9cm}@{}}
\toprule
\textbf{Fichier} & \textbf{Description} \\
\midrule
\texttt{migration-reset.sql} & Reset complet : DROP + CREATE table + RLS + 3 prospects démo \\
\texttt{migration-enrichment.sql} & Ajout colonnes enrichissement entreprise et IA \\
\texttt{migration-html-rapport.sql} & Ajout colonne \texttt{html\_rapport} \\
\texttt{migration-slug.sql} & Ajout colonne \texttt{slug} + backfill + index unique \\
\texttt{migration-notes.sql} & Ajout colonne \texttt{notes} \\
\texttt{migration-cleanup.sql} & Suppression colonnes legacy obsolètes \\
\texttt{fix-slug-duplicates.sql} & Régénération slugs en cas de doublons \\
\texttt{backfill-missing-slugs.sql} & Remplissage slugs NULL uniquement \\
\bottomrule
\end{tabular}
\end{table}

% ---------------------------------------------------------------------------
\chapter{Variables d'environnement}

\begin{table}[h]
\centering
\begin{tabular}{@{}lp{8cm}@{}}
\toprule
\textbf{Variable} & \textbf{Usage} \\
\midrule
\texttt{NEXT\_PUBLIC\_SUPABASE\_URL} & URL projet Supabase \\
\texttt{NEXT\_PUBLIC\_SUPABASE\_ANON\_KEY} & Clé publique Supabase (client) \\
\texttt{NEXT\_PUBLIC\_TAWK\_PROPERTY\_ID} & ID propriété Tawk.to \\
\texttt{NEXT\_PUBLIC\_TAWK\_WIDGET\_ID} & ID widget Tawk.to \\
\texttt{NEXT\_PUBLIC\_APP\_URL} & URL publique pour liens audit (prod) \\
\texttt{SUPABASE\_SECRET\_KEY} & Clé secrète (présente, non utilisée dans le code actuel) \\
\bottomrule
\end{tabular}
\end{table}

% ---------------------------------------------------------------------------
\chapter{Sécurité et authentification}

\section{État actuel}

\begin{itemize}[leftmargin=*]
  \item \textbf{Pas d'authentification utilisateur} dans l'application Next.js.
  \item Supabase RLS avec policies ouvertes (\texttt{USING (true)}) pour \texttt{anon} et \texttt{authenticated}.
  \item Les pages internes (dashboard, fiche prospect) ne sont pas protégées par login.
  \item Les rapports publics (\texttt{/audit/slug}) sont volontairement accessibles sans auth (partage client).
  \item Les scripts HTML des rapports sont strippés avant rendu (\texttt{stripScriptsFromHtml}).
  \item Les rapports publics sont en \texttt{noindex, nofollow}.
\end{itemize}

\section{Évolutions prévues}

\begin{itemize}[leftmargin=*]
  \item Profils multi-utilisateurs (prospecteurs avec comptes dédiés).
  \item Stats filtrées par prospecteur connecté.
  \item Authentification Supabase Auth ou NextAuth.
  \item RLS restrictif par utilisateur.
\end{itemize}

% ---------------------------------------------------------------------------
\chapter{Design system}

\section{Identité visuelle}

\begin{itemize}[leftmargin=*]
  \item \textbf{Thème} --- clair premium, inspiré Awwards.
  \item \textbf{Palette} --- indigo/violet primary, fond off-white (\texttt{oklch}), glass panels.
  \item \textbf{Typographie} --- Plus Jakarta Sans (UI), Geist Mono (chiffres/scores).
  \item \textbf{Composants} --- shadcn/ui style \texttt{base-nova}, coins arrondis 0.75--1.75rem.
  \item \textbf{Effets} --- mesh gradients, backdrop-blur, bento grids, section numbering (01, 02, 03).
  \item \textbf{Animations} --- transitions 150--200ms, hover lift sur cartes, pulse status badge.
\end{itemize}

\section{Classes utilitaires CSS}

\begin{lstlisting}
.dashboard-mesh      // Fond gradient mesh dashboard
.glass-panel         // Panneau verre (blur + ombre)
.text-gradient-brand // Titre avec degrade indigo
.card-hover-lift     // Effet elevation au survol
.prospect-page-mesh  // Mesh specifique fiche prospect
.prospect-grid-bg    // Grille de fond decorative
.section-eyebrow     // Label uppercase tracking
.bento-shine         // Reflet lumineux bento tiles
\end{lstlisting}

% ---------------------------------------------------------------------------
\chapter{Workflow commercial type}

\begin{enumerate}[leftmargin=*]
  \item \textbf{n8n} insère un prospect enrichi dans Supabase (avec ou sans slug/rapport).
  \item Le prospect apparaît dans le \textbf{pipeline} (\texttt{/}) trié par score IA.
  \item Le prospecteur consulte la \textbf{fiche} (\texttt{/prospects/id}), lit l'audit IA, prend des \textbf{notes}.
  \item Si slug manquant : clic \textbf{Générer le lien client}.
  \item Le prospecteur \textbf{copie le lien audit} et l'envoie au client (email, WhatsApp, LinkedIn).
  \item Le client ouvre \textbf{/audit/slug}, consulte son rapport brand.
  \item En restant 15s sur les faiblesses : \textbf{Mode Sniper} ouvre le chat Tawk.to avec message personnalisé.
  \item Le prospecteur \textbf{approuve} le lead depuis la carte pipeline.
  \item Le prospecteur consulte ses \textbf{stats} sur \texttt{/prospecteur}.
  \item Il copie le \textbf{script email} IA et lance sa prospection.
\end{enumerate}

% ---------------------------------------------------------------------------
\chapter{Annexes}

\section{Dépendances npm}

\subsection{Production}
\texttt{@base-ui/react}, \texttt{@supabase/supabase-js}, \texttt{class-variance-authority}, \texttt{clsx}, \texttt{lucide-react}, \texttt{next}, \texttt{react}, \texttt{react-dom}, \texttt{shadcn}, \texttt{tailwind-merge}, \texttt{tw-animate-css}.

\subsection{Développement}
\texttt{@tailwindcss/postcss}, \texttt{tailwindcss}, \texttt{typescript}, \texttt{eslint}, \texttt{eslint-config-next}.

\section{Composants non branchés}

\begin{itemize}[leftmargin=*]
  \item \texttt{prospects-table.tsx} --- vue tableau alternative (remplacée par cartes).
  \item \texttt{fast-link.tsx} --- navigation avec useTransition (remplacée par Link natif).
  \item \texttt{html-report-viewer.tsx} --- viewer iframe interne (remplacé par AuditReportEmbed).
  \item \texttt{prospect-stats.ts} --- calcul stats individuelles (préparé pour usage futur).
\end{itemize}

\section{Commandes de développement}

\begin{lstlisting}
cd target-os-dashboard
npm install
npm run dev      # http://localhost:3000
npm run build    # Build production
npm run start    # Serveur production
npm run lint     # ESLint
\end{lstlisting}

\vfill
\begin{center}
\small
\textit{Document généré pour Target OS Dashboard --- \today}\\
\textit{Version du codebase : Next.js 16.3.2 / React 19.2.8}
\end{center}

\end{document}
